% THIS DOCUMENT IS FOLLOWS THE VOLERE TEMPLATE BY Suzanne Robertson and James Robertson
% ONLY THE SECTION HEADINGS ARE PROVIDED
%
% Initial draft from https://github.com/Dieblich/volere
%
% Risks are removed because they are covered by the Hazard Analysis
\documentclass[12pt]{article}

\usepackage{booktabs}
\usepackage{tabularx}
\usepackage{hyperref}
\hypersetup{
    bookmarks=true,         % show bookmarks bar?
      colorlinks=true,      % false: boxed links; true: colored links
    linkcolor=red,          % color of internal links (change box color with linkbordercolor)
    citecolor=green,        % color of links to bibliography
    filecolor=magenta,      % color of file links
    urlcolor=cyan           % color of external links
}

\newcommand{\lips}{\textit{Insert your content here.}}

\input{../Comments.text}
\input{../Common.text}

\begin{document}

\title{Software Requirements Specification for \progname: subtitle describing software} 
\author{\authname}
\date{\today}
	
\maketitle

~\newpage

\pagenumbering{roman}

\tableofcontents

~\newpage

\section*{Revision History}

\begin{tabularx}{\textwidth}{p{3cm}p{2cm}X}
\toprule {\textbf{Date}} & {\textbf{Version}} & {\textbf{Notes}}\\
\midrule
Date 1 & 1.0 & Notes\\
Date 2 & 1.1 & Notes\\
\bottomrule
\end{tabularx}

~\\

~\newpage
\section{Purpose of the Project}
\subsection{User Business}

As AI allows the generation of code to come much more cheaply, the engineering
intent that comes behind a software system, such as its rationale and its
functional and non-functional requirements, gets lost across the quick and messy
repetitive prompt, answer, and code generation cycle. Developer teams lose the
traceability between the intent and the code/tests, which ultimately hurts
maintainability and software scalability, because through the bloated system it
is hard to find justification for why the system is built the way it is and how
the loss of control over quality and acceptance came to be.

A tool is needed that can allow traditional software engineering practices and
principles to persist and be applicable through AI agentic code and test
generation that enforces and streamlines these principles.

\subsection{Goals of the Project}

The overarching goal of \progname{} is to provide tool support that keeps human
engineers in control of engineering intent and software quality within an
AI-native (SE3.0) development workflow, where engineers define the goals,
constraints, and acceptance criteria and AI agents operationalize them. This
overarching goal is refined into the following supporting goals:

\begin{itemize}
  \item[\textbf{G1}:] Enable a developer to capture the engineering intent
  behind a change from natural-language interaction with an AI coding assistant,
  and require explicit developer approval of that intent before any code is
  generated.
  \item[\textbf{G2}:] Maintain a versioned intent model which consists of the following: the software's
  requirements, quality constraints, and design rationale. This will be a durable artifact that may be in a read-only 
  setting for the user so requirement linage is not tampered with and it serves as a source of truth for the software application journey and evolution. 
  \item[\textbf{G3}:] Maintain traceability links from each element of the
  intent model to the code and tests that implement it, each with a recorded
  justification.
  \item[\textbf{G4}:] Verify changes proposed or generated by an AI coding
  assistant against the intent model, and report whether they satisfy the
  recorded requirements and constraints.
  \item[\textbf{G5}:] Keep the developer in control of the system's engineering
  intent, so that every change to the intent model is reviewed and approved by a
  human before it takes effect (always have a human-in-the-loop). 
\end{itemize}
\section{Stakeholders}
\subsection{Client}

The client for \progname{} is Dr. Istvan David, the project supervisor and a
faculty member in the Department of Computing and Software (CAS) at McMaster
University. The client defines the vision of the project, sets the
acceptance criteria for the prototype, and provides ongoing feedback and direction. 

\subsection{Customer}

The customers for \progname{} are software engineering development teams that
build and evolve software with the help of AI coding assistants. These teams
need to manage their requirements, quality constraints, and the evolution of
their application through code that is written and changed across multiple
developers. For these customers, engineering intent is easily lost as work is
spread across many contributors and there custom agentic systems, so they need a way
to keep engineering principles connected to the code and test cases that implement them. 
\progname{} is intended to serve these teams by keeping that intent explicit, traceable, and under human control.

\subsection{Other Stakeholders}

Table~\ref{TblOtherStakeholders} lists stakeholders who are not users of
\progname{} but who influence or are affected by it.

\begin{table}[h!]
\centering
\caption{Other Stakeholders}
\label{TblOtherStakeholders}
\begin{tabularx}{\textwidth}{p{4.2cm}X}
\toprule
\textbf{Stakeholder} & \textbf{Interest in \progname{}} \\
\midrule
Development team (Team 8) & Designs, builds, and tests the product. Owns the
requirements and documentation. \\
Course instructor and TAs & Evaluate the documentation, demonstrations, and
final product against the course rubrics. \\
AI coding assistant vendors & Provide the assistants \progname{} integrates
with, such as Claude Code, Codex, and Cursor. Changes to their MCP support can
break the integration. \\
LLM and tracing service providers & Provide the LLM used for extraction and
verification, and LangSmith for tracing. They receive project data sent by the
product. \\
IDE vendors and extension marketplaces & Define the extension APIs and
publishing rules the product must follow, such as the VS Code Marketplace. \\
Open-source license holders & Own the libraries the product uses. Their
licenses must be respected. \\
End users of software built with \progname{} & Indirectly affected. They rely
on software whose quality depends on the intent captured by the product. \\
\bottomrule
\end{tabularx}
\end{table}

\subsection{Hands-On Users of the Project}

\progname{} has three categories of hands-on users.

\begin{itemize}
  \item \textbf{AI-assisted developer.} Writes code with an AI coding assistant
  inside an IDE. Submits requests, reviews and approves proposed requirements,
  and reads verification results.
  \begin{itemize}
    \item Subject matter experience: knows the software being built.
    Familiar with requirements and design concepts, but may not write formal
    requirements.
    \item Technological experience: comfortable with an IDE, Git, and at least
    one AI coding assistant.
    \item Education: undergraduate degree in software engineering, computer
    science, or equivalent experience.
    \item Attitude toward technology: uses AI tools daily and wants to keep
    their speed. Will avoid steps that feel slow or repetitive.
  \end{itemize}

  \item \textbf{Technical lead or reviewer.} Reviews changes made by other
  developers. Uses the intent model, traceability links, and verification
  results to judge whether a change meets its requirements.
  \begin{itemize}
    \item Subject matter experience: high. Understands the system design and
    its quality priorities.
    \item Technological experience: high. Uses Git, GitHub pull requests, and
    code review tools.
    \item Attitude toward technology: values traceability and control over
    speed. Sceptical of AI-generated code.
  \end{itemize}

  \item \textbf{New team member.} Joins a project that already uses
  \progname{}. Reads the intent model to learn why the code is structured the
  way it is.
  \begin{itemize}
    \item Subject matter experience: low for the project being joined.
    \item Technological experience: comfortable with an IDE and Git.
    \item Attitude toward technology: relies on documentation and existing
    artifacts to learn the codebase.
  \end{itemize}
\end{itemize}

\subsection{Personas}

\textbf{Priya, AI-assisted developer.} Priya is a full-stack developer at a
small software company. She uses Claude Code in VS Code for most new features.
Her team has lost track of why parts of the code were written, because the
reasoning only existed in old chat sessions. Priya wants her requirements
recorded without slowing down. She will use \progname{} if reviewing a proposal
takes less time than explaining the feature again later.

\textbf{Daniel, technical lead.} Daniel reviews pull requests for a team of six
developers who all use AI assistants. He cannot read every line of generated
code, so he needs to know which requirements a change is meant to satisfy and
whether it does. He uses the intent model and verification results during code
review.

\textbf{Sam, new team member.} Sam joined Daniel's team last month. When Sam
finds code that looks unnecessary, Sam checks the intent model to find the
requirement and rationale behind it before changing anything.

\subsection{Priorities Assigned to Users}

\begin{itemize}
  \item \textbf{Key users:} AI-assisted developers. The product cannot meet its
  goals unless they capture and approve intent (G1, G5).
  \item \textbf{Secondary users:} technical leads, reviewers, and new team
  members. They read the intent model, traceability links, and verification
  results, but do not create most of the content.
  \item \textbf{Unimportant users:} non-technical stakeholders, such as product
  managers, who may read the intent model. The product does not provide
  features specifically for them.
\end{itemize}

When the needs of key and secondary users conflict, the needs of key users take
priority.

\subsection{User Participation}

\begin{itemize}
  \item \textbf{Client.} Dr.~David meets with the team weekly to review
  requirements and progress, and sets the acceptance criteria for the
  prototype.
  \item \textbf{Development team.} Team members are AI-assisted developers
  themselves and represent the key user group when writing and reviewing
  requirements.
  \item \textbf{Representative users.} Developers outside the team will try
  the review and approval workflow during validation. Details will be given in
  the V\&V Plan.
  \item \textbf{Demonstrations.} Feedback from the Proof of Concept and
  Revision~0 demonstrations will be used to revise the requirements.
\end{itemize}

\subsection{Maintenance Users and Service Technicians}

\progname{} is software only, so there are no service technicians.

During the course, the development team maintains the product. After the
course, maintenance responsibility will be agreed with the client.

Maintainers need:
\begin{itemize}
  \item documentation of each major component and interface
  (Maintenance Requirements);
  \item diagnostic logs for significant failures (Supportability
  Requirements);
  \item isolated integrations so new AI coding assistants can be added
  (Adaptability Requirements).
\end{itemize}

Developers who use \progname{} update it through their IDE's extension
marketplace and rely on release notes to learn what changed (REL1).

\section{Mandated Constraints}
\subsection{Solution Constraints}

\medskip
\noindent\textbf{Description:} The product shall integrate with and run inside
multiple integrated development environments (IDEs). VS Code and Cursor are the
primary targets, and PyCharm is a secondary target if time permits. \\
\textbf{Rationale:} Developers work in different IDEs, so the tool must meet them
in the environment they already use rather than forcing them to adopt a new one. \\
\textbf{Fit Criterion:} The core intent-capture and approval workflow shall be
verified to run inside VS Code and Cursor, and support for PyCharm shall be
attempted if the schedule allows.

\medskip
\noindent\textbf{Description:} The product shall function with multiple agentic
coding assistants. Claude Code is the primary target, and Codex and Grok or
Cursor AI are secondary targets if time permits. \\
\textbf{Rationale:} Teams rely on different AI coding assistants, and the tool
must capture intent and constrain generated code regardless of which assistant
produces it. \\
\textbf{Fit Criterion:} The product shall be verified to capture intent and gate
code generation when driven by Claude Code, and compatibility with Codex and
Grok or Cursor AI shall be attempted if the schedule allows.

\subsection{Implementation Environment of the Current System}

In the current design, \progname{} runs locally on the developer's machine as an
IDE extension together with the IDE terminal running a AI coding assistant or 
an IDE plugin tool. The extension and server coordinate through
files in the developer's workspace, and requirement extraction verification may call an
external large language model service over the network. The product shall
function on the operating systems that host the supported IDEs, namely Windows,
macOS, and Linux.

\subsection{Partner or Collaborative Applications}

\progname{} is designed to work alongside the developer's existing AI coding
assistant and host IDE rather than replace them, and it integrates with the
project's version control system (Git and GitHub) to version the intent model.
These are the applications the product must cooperate with, and there are no
other partner or collaborative constraints at this time.

\subsection{Off-the-Shelf Software}

The following off-the-shelf software and frameworks will be utilized:
\begin{itemize}
  \item Python - V3.14
  \item TypeScript - V7.0
  \item LangGraph
  \item LangSmith
\end{itemize}

\subsection{Anticipated Workplace Environment}

The anticipated workplace environment is a software developer working at a
desktop or laptop computer, in an office or remote setting, using an IDE paired
with an AI coding assistant. The product is used wherever the developer already
carries out their development work, and it imposes no additional physical
environment requirements.

\subsection{Schedule Constraints}

As set out in the course schedule, the project must meet the following
deliverable deadlines.

\begin{table}[h!]
\centering
\caption{Schedule Constraints}
\label{TblSchedule}
\begin{tabular}{l p{8.5cm}}
\toprule
\textbf{Date} & \textbf{Deliverable} \\
\midrule
Oct 19, 2026       & Requirements Document (SRS) and Hazard Analysis \\
Nov 2, 2026        & Verification and Validation Plan \\
Nov 30 to Dec 3, 2026 & Proof of Concept Demonstration \\
Jan 18, 2027        & Design Document \\
Feb 3 to Feb 12, 2027 & Revision 0 Demonstration \\
Mar 8, 2027        & Verification and Validation Report \\
Mar 18 to Mar 20, 2027 & Final Demonstration (Revision 1) \\
Apr 5, 2027       & Final Documentation \\
Apr 6, 2027         & Capstone EXPO \\
\bottomrule
\end{tabular}
\end{table}

\subsection{Budget Constraints}

The project has no monetary budget. Development currently relies on free and
open-source tooling. 
Any future cost would be limited to optional paid software services, for example a LangSmith subscription
for tracing and evaluation, a hosted MCP prompt-versioning tool, or a managed
database for storing the intent model. If any purchase is required, it would be discussed with and approved by the
supervisor beforehand.

\subsection{Enterprise Constraints}

\progname{} is a local, single-developer tool, and it does not currently impose
enterprise-level constraints such as centralized authentication or multi-tenant
access control. The main enterprise constraint is that the tool must fit into
existing software engineering workflows and version control practices without
requiring teams to adopt new infrastructure. 

\section{Naming Conventions and Terminology}
\subsection{Glossary of All Terms, Including Acronyms, Used by Stakeholders
involved in the Project}

\begin{description}
  \item[AI coding assistant] An AI tool that writes or changes code from
  natural-language requests, such as Claude Code, Codex, or Cursor AI.
  \item[Agentic] Describes an AI system that carries out multi-step tasks,
  such as editing files and running commands, with limited human input.
  \item[Approval] The developer's decision to accept a proposal. Only approved
  requirements enter the intent model and guide code generation.
  \item[Design constraint] A restriction on how the software may be designed or
  built, such as a required framework or architecture.
  \item[Design rationale] The reason a requirement or design decision was made.
  \item[Engineering artifact] Any item recorded by \progname{}, such as a
  requirement, constraint, design decision, or traceability link.
  \item[Engineering intent] What the developer wants the software to do and
  why, including requirements, quality priorities, constraints, and rationale.
  \item[Fit criterion] A measurable condition used to decide whether a
  requirement has been met.
  \item[Functional requirement] A behaviour the system must perform.
  \item[Git] A version control system for tracking changes to files.
  \item[GitHub] A hosting service for Git repositories.
  \item[Host IDE] The IDE in which \progname{} runs, such as VS Code, Cursor,
  or PyCharm.
  \item[Human-in-the-loop] A workflow in which a person must review and approve
  an AI result before it takes effect.
  \item[IDE] Integrated Development Environment. An application for writing,
  running, and debugging code.
  \item[Intent model] The versioned, structured record of a project's approved
  requirements, quality constraints, design rationale, and traceability links.
  It is the source of truth for engineering intent.
  \item[LangGraph] A framework for building LLM workflows. \progname{} uses it
  to run extraction with human approval steps.
  \item[LangSmith] A service for tracing and evaluating LLM workflows.
  \item[LLM] Large Language Model. An AI model that generates text, used by
  \progname{} to extract and verify requirements.
  \item[MCP] Model Context Protocol. An open protocol that lets AI coding
  assistants call external tools. \progname{} uses it to communicate with the
  assistant.
  \item[Non-functional requirement] A quality the system must have, such as
  performance, security, or usability.
  \item[Proposal] A set of requirements and constraints extracted from a
  developer's request and awaiting approval, editing, or rejection. Also called
  a draft proposal.
  \item[Quality constraint] A required level of a quality attribute, such as
  maximum response time.
  \item[Requirement lineage] The recorded history of how a requirement was
  created and changed.
  \item[Schema] The defined structure that intent model data must follow.
  \item[SE3.0] Software Engineering 3.0. A development approach in which
  engineers define goals, constraints, and acceptance criteria, and AI agents
  carry out the implementation.
  \item[SRS] Software Requirements Specification. This document.
  \item[Traceability link] A recorded, justified connection between an
  element of the intent model and the code or tests that implement it.
  \item[Verification] Checking a code change against the intent model to
  decide whether it satisfies the recorded requirements and constraints.
  \item[VSIX] The package format for VS Code extensions.
  \item[Workspace] The project folder open in the developer's IDE.
\end{description}

\section{Relevant Facts And Assumptions}
\subsection{Relevant Facts}
\lips
\subsection{Business Rules}
\lips
\subsection{Assumptions}
\lips

\section{The Scope of the Work}
\subsection{The Current Situation}
\lips
\subsection{The Context of the Work}
\lips
\subsection{Work Partitioning}
\lips
\subsection{Specifying a Business Use Case (BUC)}
\lips

\section{Business Data Model and Data Dictionary}
\subsection{Business Data Model}
\lips
\subsection{Data Dictionary}
\lips

\section{The Scope of the Product}
\subsection{Product Boundary}
\lips
\subsection{Product Use Case Table}
\lips
\subsection{Individual Product Use Cases (PUC's)}
\lips

\section{Functional Requirements}

The functional requirements describe the behaviour \progname{} must provide to
meet the goals of the project.
 A functional requirement is a baseline acceptance criteria of what the system should be able to do. 

\subsection{Intent Capture}

\medskip
\noindent\textbf{FR1.} The system shall accept a developer's natural-language
feature request that is submitted through the AI coding assistant that is running anywhere in there IDE. 

\medskip
\noindent\textbf{FR2.} The system shall extract a schema specified structured set of proposed
functional requirements, non-functional requirement, quality constraints, and design constraints from the developer's request.

\medskip
\noindent\textbf{FR3.} The system shall save the extracted requirements and
quality constraints as a draft proposal and make that proposal available for the
developer to review and then accept, reject, or follow up/edit. 

\subsection{Requirement Review and Approval}

\medskip
\noindent\textbf{FR4.} The system shall present the extracted functional
requirements and quality constraints to the developer for review before any
corresponding code is generated.

\medskip
\noindent\textbf{FR5.} The system shall allow the developer to approve, edit, or reject
a proposal.

\medskip
\noindent\textbf{FR6.} The system shall allow the developer to remove or amend
individual requirements or quality constraints within a proposal while it is
still pending.

\subsection{Code Generation Control}

\medskip
\noindent\textbf{FR7.} The system shall prevent the AI coding assistant from
generating code for a request until the developer has approved or rejected the
corresponding proposal.

\medskip
\noindent\textbf{FR8.} The system shall communicate the developer's decision
back to the AI coding assistant, and the tools internal database which manages the requirements and constraints, instructing it to implement exactly the
approved requirements when the proposal is approved, or to not implement and
instead ask how to revise the request when the proposal is rejected if a edit is not suggested when rejecting the initial proposed requirment/constraint. 

\subsection{Intent Model and Traceability}

\medskip
\noindent\textbf{FR9.} The system shall record the approved requirements,
quality constraints, and design rationale in a versioned intent model with a structued data schema, and shall
maintain traceability links from these elements to the code and tests that
implement and satisfy respectively.

\subsection{Change Verification}

\medskip
\noindent\textbf{FR10.} The system shall check changes proposed or generated by
the AI coding assistant against the intent model and report whether they satisfy
the recorded requirements and constraints.

\subsection{Design Rationale}

\medskip
\noindent\textbf{FR11.} The system shall record, for each approved requirement
and design decision, the rationale behind it and the originating developer
request that led to it, so that it is clear why a certain requirement exists. 

\subsection{Constrained Generation}

\medskip
\noindent\textbf{FR13.} The system shall compile the approved requirements and
quality constraints into executable test, that the AI-generated code assist must satisfy.

\section{Look and Feel Requirements}
\subsection{Appearance Requirements}

\noindent\textbf{AR1.} The system shall present its interface in a way that integrates cleanly within the host IDE and does not unnecessarily disrupt the developer's existing workspace.

\noindent\textbf{AR2.} The interface shall use a consistent visual structure, comprised of readable fonts, colours, spacing, and logical grouping of related information, to keep engineering artifacts and workflow information easy to understand.

\subsection{Style Requirements}

\noindent\textbf{SR1.} The system shall maintain a professional visual style that is consistent with the appearance and interaction patterns of the host IDE.

\noindent\textbf{SR2.} The system shall maintain a consistent visual style across its views, prompts, artifact displays, and workflow states.

\section{Usability and Humanity Requirements}
\subsection{Ease of Use Requirements}

\noindent\textbf{EUR1.} The system shall provide a simple and intuitive workflow for submitting development requests, reviewing proposed engineering artifacts, providing further detail / specification, and approving, rejecting, or modifying them.

\noindent\textbf{EUR2.} The system shall clearly indicate the actions available to the user at each stage of the workflow so that approval, rejection, revision, and continuation steps are easy to understand.

\noindent\textbf{EUR3.} The system shall make traceability information easy to locate by allowing users to identify the code and tests associated with a selected engineering artifact.

\subsection{Personalization and Internationalization Requirements}

\noindent\textbf{PIR1.} The system shall use clear and culturally neutral language, symbols, and interaction patterns so that the core workflow is understandable to users from different backgrounds.

\subsection{Learning Requirements}

\noindent\textbf{LR1.} The system shall be easy to learn by using a simple interaction based workflow consisting of a clear exchange of information between the developer and the system.

\noindent\textbf{LR2.} The system shall require minimal instruction for a new user to understand how to provide information and respond to the system's requests for clarification or approval.

\subsection{Understandability and Politeness Requirements}

\noindent\textbf{UPR1.} The system shall communicate with developers using clear software engineering terminology and shall not require users to understand internal implementation details of the system.

\noindent\textbf{UPR2.} The system shall avoid requesting information from the developer when that information can already be obtained from the project workspace or existing engineering artifacts.

\subsection{Accessibility Requirements}

\noindent\textbf{AR1.} The system shall operate within the host IDE and shall not require users to learn or operate a separate standalone interface.

\noindent\textbf{AR2.} The system shall use standard IDE interaction patterns and controls where possible so that existing accessibility features provided by the host IDE can continue to be used.

\section{Performance Requirements}
\subsection{Speed and Latency Requirements}

\noindent\textbf{PR1.} The system shall process developer requests and produce engineering artifact proposals within a reasonable amount of time relative to a typical AI coding assistant interaction, such that the additional IDE3 processing does not significantly disrupt the developer's workflow.

\noindent\textbf{PR2.} The system shall provide immediate user feedback when a long running operation begins, so that the developer is aware that the request is being processed.

\subsection{Safety-Critical Requirements}

\noindent\textbf{SCR1.} The system shall require explicit human approval of proposed engineering artifacts before they are used to guide code generation or modifications.

\subsection{Precision or Accuracy Requirements}

\noindent\textbf{PAR1.} When a developer rejects or modifies a proposed engineering artifact, the system shall incorporate the provided feedback into subsequent revisions until an acceptable set of artifacts is produced or the workflow is terminated by the developer.

\subsection{Robustness or Fault-Tolerance Requirements}

\noindent\textbf{RFR1.} The system shall preserve approved engineering artifacts and their relationships if an active workflow fails or is interrupted.

\noindent\textbf{RFR2.} A failure in an external service, coding agent, or model interaction shall not corrupt previously persisted project data.

\noindent\textbf{RFR3.} Where possible, the system shall allow a developer to resume an interrupted workflow without repeating previously completed and approved steps.

\subsection{Capacity Requirements}

\noindent\textbf{CR1.} The system shall be capable of maintaining engineering artifacts, their relationships, and traceability information for software projects of a size representative of a medium sized development project.

\noindent\textbf{CR2.} The system shall support processing development requests that produce multiple related engineering artifacts without requiring the user to manually divide the request into separate interactions.

\subsection{Scalability or Extensibility Requirements}

\noindent\textbf{SER1.} The system shall maintain acceptable performance as the size of the project, number of engineering artifacts, and number of traceability relationships increase.

\subsection{Scalability or Extensibility Requirements}

\noindent\textbf{SER1.} The system shall maintain acceptable performance as the size of the project, number of engineering artifacts, and number of traceability relationships increase.

\noindent\textbf{SER2.} The system shall be structured so that support for additional coding agents, engineering artifact types, and analysis methods can be added without significant restructuring of the core system.

\subsection{Longevity Requirements}

\noindent\textbf{LR1.} The system shall be designed to remain adaptable as AI models, coding agents, engineering practices, and forms of engineering input evolve.

\section{Operational and Environmental Requirements}
\subsection{Expected Physical Environment}

Not applicable. 

\subsection{Wider Environment Requirements}

Not applicable. \progname{} operates entirely within the developer's local IDE
and workspace and imposes no requirements arising from a wider physical, environment. 

\subsection{Requirements for Interfacing with Adjacent Systems}

\medskip
\noindent\textbf{IAS1.} The system shall interface with each supported host IDE
through that IDE's official extension API, such as the VS Code Extension API. 

\medskip
\noindent\textbf{IAS2.} The system shall communicate with the AI coding
assistant through the Model Context Protocol (MCP), so that any assistant tool that
supports MCP can drive the intent-capture, code, and test workflow. 

\medskip
\noindent\textbf{IAS3.} The system shall interface with the project's version
control system, such as Git, using standard Git operations to version the
intent model.

\subsection{Productization Requirements}

\medskip
\noindent\textbf{PRD1.} The system shall be packaged as an installable IDE
extension in the outlined format for each supported IDE, such as a VSIX package
for VS Code, so that it can be installed by the user. 

\medskip
\noindent\textbf{PRD2.} The system shall satisfy the publishing and security
requirements of the relevant extension marketplace, such as the Visual Studio
Code Marketplace, before it is accepted to be discoverable and installable from the marketplace. 

\medskip
\noindent\textbf{PRD3.} The system shall bundle its runtime
dependencies, so that installation does not require the user to set up those dependencies manually.

\subsection{Release Requirements}

\medskip
\noindent\textbf{REL1.} Each release shall be accompanied by release notes that
describe the new features added and the bugs or errors patched since the previous
version.

\medskip
\noindent\textbf{REL2.} The system shall preserve compatibility of the intent
model format across releases, or provide a migration path, so that new releases
do not introduce any breaking changes to existing intent models. 

\section{Maintainability and Support Requirements}
\subsection{Maintenance Requirements}

\noindent\textbf{MR1.} The system shall organize major functionality into modular components with clearly defined responsibilities and interfaces to ensure that changes to one component have minimal impact on unrelated components.

\noindent\textbf{MR2.} Dependencies on external libraries, frameworks, models, and services shall be isolated where possible so that they can be upgraded or replaced without requiring major changes throughout the system.

\noindent\textbf{MR3.} Integrations with external coding agents shall be separated from the core intent management and traceability functionality so that additional agents can be supported without significant restructuring of the system.

\noindent\textbf{MR4.} Major components and interfaces shall have accompanying documentation describing their purpose, dependencies, and expected usage to support future maintenance and modification.

\subsection{Supportability Requirements}

\noindent\textbf{SR1.} The system shall handle failures gracefully and provide clear error messages that describe the cause of the issue where it can be identified, such as failures involving external coding agents, model services, or other system components.

\noindent\textbf{SR2.} The system shall record diagnostic logs for significant failures and system interactions so that recurring issues can be identified and investigated during development and maintenance.

\subsection{Adaptability Requirements}

\noindent\textbf{AR1.} The system shall isolate integrations with external coding agents so that support for additional agents can be added without significant changes to the core intent management and traceability functionality.

\noindent\textbf{AR2.} The system shall be designed so that additional forms of engineering input (e.g. Petri nets) can be incorporated without requiring significant changes to the core intent model or processing workflow.

\section{Security Requirements}

\subsection{Access Requirements}

\medskip
\noindent\textbf{ACR1.} Only the local developer who owns the code directory workspace shall be
able to review, approve, reject, or edit a proposal, and the AI coding assistant
shall never be able to approve, edit, or reject potential proposals. 

\medskip
\noindent\textbf{ACR2.} The system shall restrict read and write access to the
intent model to the user who owns the developer's workspace. 

\medskip
\noindent\textbf{ACR3.} The system shall derive access to the intent model from
the access a user already has to the project's codebase, so that a user with
write access to the repository has write access to the intent model. The
system shall inherit the permissions enforced by the project's repository rather
than maintain a separate access-control scheme of its own.

\subsection{Integrity Requirements}

\medskip
\noindent\textbf{IR1 (Atomicity).} The system shall apply each change 
to the intent model as a single atomic operation, so that the change
either takes effect in full or has no effect at all, this is so a reader never observes
a partially applied or half-written change. An approval or rejection that does
not complete, shall leave the intent model in its most recent whole complete state.

\medskip
\noindent\textbf{IR2 (Consistency).} The system shall ensure that every
committed change leaves the intent model in a valid state that satisfies the schema rules. 

\medskip
\noindent\textbf{IR3 (Isolation).} The system shall isolate concurrent
operations on the intent model so that they do not interfere
with one another. 

\medskip
\noindent\textbf{IR4 (Durability).} The system shall ensure that once a change is
committed it is persisted durably.
Committed intent(s) shall be versioned so that the state of the intent
model can be recovered and its history reviewed so that nothing is silently lost or unrecoverable.

\subsection{Privacy Requirements}

\medskip
\noindent\textbf{PRR1.} The system shall store the intent model only within the developer's workspace, 
and shall not share them with third party services until and unless the developer has accepted risks for sharing it with developer tools such as version control like Github.

\subsection{Audit Requirements}

\medskip
\noindent\textbf{AUR1.} The system shall record, for each requirement/constraint decision,
the decision outcome, and the time the decision was made, and the time a pre-existing decision was modified. 

\medskip
\noindent\textbf{AUR2.} The system shall maintain a reviewable history and lineage of the 
changes to the intent model, including what changed, when it changed, and the
rationale for the change.

\medskip
\noindent\textbf{AUR3.} The system shall record each verification of an
AI-proposed or AI-generated change against the intent model.

\medskip
\noindent\textbf{AUR4.} For any future multi-user deployment, the system shall
attribute each decision and each change to the user who performed it.

\subsection{Immunity Requirements}

\medskip
\noindent\textbf{IMR1.} The system shall validate and bound the size of each
incoming request, and shall reject empty, oversized, or malformed input, to
protect against malformed-input and resource-exhaustion attacks.

\medskip
\noindent\textbf{IMR2.} The system shall fail safely when the server, review
panel, or extraction service becomes unreachable, the system should fall to a state in the application that
does not generates code compared to an unsafe and unpredictable state. 

\medskip
\noindent\textbf{IMR3.} The system shall keep its third-party dependencies up to date, and each dependency upgrade shall be accompanied by the software changes that reflect the reason for that upgrade.
Version bump alone is not enough. 

\section{Cultural Requirements}?
\subsection{Cultural Requirements}
\noindent\textbf{CR1.} The system shall communicate via Canadian English. 
\medskip
\noindent\textbf{CR2.} The system shall avoid offens ive language. If the system begins to accept image (e.g., for architectural review purposes), then no image pertaining to or related in any form to religion, political views, or minority groups will be accepted.
\lips

\section{Compliance Requirements}
\subsection{Legal Requirements}
\medskip
\noindent\textbf{LR1.} The system shall not infringe any copyright or license through the utilization of assets such as public code, system architectures, or open-source tools.
If any open-source or third-party resources are used, then the system shall explicitly cite those resources, comply with the terms and conditions of each resource's license, and request permission for any license that requires it before the system is distributed publicly.
\subsection{Standards Compliance Requirements}
\medskip
\noindent\textbf{SCR1.} The system shall conform to the Model Context Protocol (MCP) specification for any AI tool it integrates with.

\medskip
\noindent\textbf{SCR2.} The system shall comply with the VS Code Extension API guidelines and their equivalents for the other supported IDEs, such as Cursor and PyCharm.

\medskip
\noindent\textbf{SCR3.} The system shall persist the intent model using its defined schema, in a documented standard format, such as JSON or MD or any other developed model language, that is consistent across all instances of the tool.

\medskip
\noindent\textbf{SCR4.} The system shall conform to the host IDE's accessibility guidelines for any interface exposed to the user.

\section{Open Issues}
There are currently no open issues to be discussed. 
\lips

\section{Off-the-Shelf Solutions}
\subsection{Ready-Made Products}
\lips
\subsection{Reusable Components}
\lips
\subsection{Products That Can Be Copied}
\lips

\section{New Problems}
\subsection{Effects on the Current Environment}
\lips
\subsection{Effects on the Installed Systems}
\lips
\subsection{Potential User Problems}
\lips
\subsection{Limitations in the Anticipated Implementation Environment That May
Inhibit the New Product}
\lips
\subsection{Follow-Up Problems}
\lips

\section{Tasks}
\subsection{Project Planning}
\lips
\subsection{Planning of the Development Phases}
\lips

\section{Migration to the New Product}
\subsection{Requirements for Migration to the New Product}
\lips
\subsection{Data That Has to be Modified or Translated for the New System}
\lips

\section{Costs}
\lips
\section{User Documentation and Training}
\subsection{User Documentation Requirements}
\lips
\subsection{Training Requirements}
\lips

\section{Waiting Room}
\lips

\section{Ideas for Solution}
\lips

\newpage{}
\section{Appendix --- Generative AI Use Disclosure}

\wss{ChatGPT (version 5) was used to brainstorm ideas for the template for this disclose section.}

\subsection{AI Tools Used}

\wss{Give the name of the AI tool(s) used and the version.}

\subsection{How and Where Used}

\wss{Explain what the tool(s) were used for.  Examples include: brainstorming,
explaining a technical concept, reviewing the document, generating or modifying
text, generating or modifying diagrams, identifying errors or omissions or
inconsistencies.}

\wss{For each of the uses, identify which parts of the document were affected.}

\wss{You do not need to report every interaction, but you should disclose the 
uses that materially contributed to your work}

\subsection{Verification of AI-Generated Content}

\wss{\begin{itemize}
    \item Describe how the team checked AI-generated or AI-assisted material
    before including it in the document.
    \item In particular, verify factual claims, technical assertions,
    requirements, calculations, references, and other information for which an
    AI system may produce incorrect or fabricated results.
    \item \textbf{The team is responsible for the accuracy, completeness,
consistency, and appropriateness of the submitted document, regardless of
whether material was generated by a human or an AI system.}
\end{itemize}}

\newpage{}

\section{Appendix --- Reflection}

\input{../Reflection.text}

\input{../SRS_Reflection.text}

\end{document}